\section{Discussion}
\label{sec:discussion}

\subsection{What is fixed by the continuum theory fragment}

At the continuum level, the theory specifies:
(i) the substrate dynamics via an isotropic hyperelastic action \eqref{eq:action},
(ii) the existence of an ordered narrowband sector that supports a controlled multiscale envelope description,
and (iii) the geometric connection structure (Berry/Wilczek--Zee) induced by polarization transport.
Given these commitments, gauge degrees of freedom are not introduced as independent fundamental fields; they arise as
connection data of the ordered carrier sector.

\subsection{de~Broglie--Einstein ``harmony of phases''}

The classical de~Broglie picture emphasizes a relation between an internal ``clock'' (Compton-scale oscillation)
and an externally observed matter-wave phase.
In the present framework this relation is realized naturally: a localized soliton can support an internal carrier-scale
oscillation propagating at the microscopic wave speed $c$ in the effective substrate/branch metric, while the soliton
center-of-mass follows a timelike trajectory in the emergent kinematics.
The slow envelope phase then encodes the externally observed matter-wave behavior, while the internal carrier provides
the phase reference that enables Berry/Wilczek--Zee holonomy.

\subsection{Quantum phenomenology without fundamental probability}

By construction, the underlying dynamics is deterministic.
The probabilistic character of quantum measurements is attributed to the interaction of localized solitons with
wave packets that are constrained by Fourier localization (band-limited narrowband carriers cannot be localized
arbitrarily sharply), together with threshold-like nonlinear exchange events in which small differences in phase and
amplitude can decide whether energy is transferred.
In this view, apparent point-like interactions reflect the structure of wave packets and nonlinear exchange,
not the existence of fundamental point particles.

\subsection{Limitations and open problems}

Several gaps remain before the theory fragment could be regarded as a closed replacement for established field theories:
\begin{itemize}
\item identifying constitutive laws and parameter regimes that simultaneously support (a) a stable ordered narrowband
sector and (b) robust localized solitons,
\item deriving the precise effective dynamics of curvature fluctuations of the Berry/Wilczek--Zee connection
(including whether a Maxwell/Yang--Mills-type term dominates in the relevant regimes),
\item clarifying the quantitative mapping between far-field holonomy labels and electromagnetic phenomenology,
\item understanding how multi-particle states and apparent entanglement-like correlations arise in a deterministic
wave medium.
\end{itemize}
The present goal is to place the theory on a consistent continuum footing so that these problems can be attacked
without ambiguity about the underlying substrate model.
